\documentclass[10pt,UTF8]{article}

\usepackage{geometry}

\usepackage{ctex}

\geometry{a4paper,scale=0.9}
\ctexset{today=old}

\usepackage[bookmarksnumbered=true]{hyperref}  % 生成大纲



\usepackage{times}  % 设置正文字体为 Adobe Times Roman

% \usepackage{makecell}
% \usepackage{multirow} % 表格多行合并

% \usepackage{graphicx} %插入图片的宏包
% \usepackage{subfigure}
% \usepackage{float} %设置图片浮动位置的宏包

\usepackage{amsmath}
\usepackage{amssymb}
\usepackage{mathtools}

\usepackage{physics}

% \usepackage{siunitx}
% % 关闭警告
% \ExplSyntaxOn
% \msg_redirect_name:nnn { siunitx } { physics-pkg } { none }
% \ExplSyntaxOff

\allowdisplaybreaks[3]  % 允许换页


%%%%%%%%%%%%%%%%%%%%%%%% tcolorbox %%%%%%%%%%%%%%%%%%%%%%%%%%
\usepackage[most]{tcolorbox}

\newenvironment{definition box}
{\begin{tcolorbox}[colback=blue!5!white,colframe=blue!75!black,parbox=false,before upper=\par,before lower=\par]}
{\end{tcolorbox}}

\newenvironment{theorem box}
{\begin{tcolorbox}[colback=yellow!5!white,colframe=yellow!75!black,parbox=false,before upper=\par,before lower=\par]}
{\end{tcolorbox}}
%%%%%%%%%%%%%%%%%%%%%%%%%%%%%%%%%%%%%%%%%%%%%%%%%%%%%%%%%%%%
            
%%%%%%%%%%%%%%%%%%%%% amsthm %%%%%%%%%%%%%%%%
\usepackage{amsthm}

\theoremstyle{definition}
\newtheorem{definition inner}{Definition}[section]
\newenvironment{definition}[1][]
{\begin{definition box}\begin{definition inner}[#1]\hfill\par}
{\end{definition inner}\end{definition box}}

\theoremstyle{plain}
\newtheorem{theorem inner}{Theorem}[section]
\newenvironment{theorem}[1][]
{\begin{theorem box}\begin{theorem inner}[#1]\hfill\par}
{\end{theorem inner}\end{theorem box}}

\theoremstyle{definition}
\newtheorem*{proof inner}{\textit{Proof}}
\renewenvironment{proof}[1][]
{\begin{proof inner}[#1]\hfill\par}
{\qed\end{proof inner}}

\theoremstyle{remark}
\newtheorem*{remark}{\textit{Remark}}
%%%%%%%%%%%%%%%%%%%%%%%%%%%%%%%%%%%%%%%%%%%%%%


\newcommand{\mycases}[2][0.8]{
    \left\{\vphantom{\scalebox{#1}{
        $\begin{matrix}#2\end{matrix}$
    }}\right.\!\!\!\!
    \begin{array}{l@{\quad}l}#2\end{array}
}

\newcommand{\ju}{\mathrm{j}}  % 虚数单位
\newcommand{\e}{\mathrm{e}{\mkern 1 mu}}

\newcommand{\Ev}{\,\mathcal{E}\ensuremath{v}\qty}
\newcommand{\Od}{\,\mathcal{O}\ensuremath{d}\qty}

\renewcommand{\Re}{\,\mathcal{R}\ensuremath{e}\qty}
\renewcommand{\Im}{\,\mathcal{I}\ensuremath{m}\qty}

\newcommand{\sinc}{\operatorname{sinc}} 
\newcommand{\rect}{{\mkern 1.5 mu}\operatorname{rect}\!\qty} % 矩形函数


\begin{document}

\part*{The Discrete-Time Fourier Transform}

\section{Representation of DTFT}


\begin{definition}[The Fourier Transform Pair of a DT Aperiodic Signal]
    If an aperiodic signal $x[n]$ has a Fourier transform, then
    \begin{align*}
        x[n] &=
        \frac{1}{2\pi}\int_{2\pi}
        X(\e^{\ju\omega})\e^{\ju\omega n}\dd{\omega}
        \tag{The Synthesis Equation} \\
        X(\e^{\ju\omega}) &=
        \sum_{n=-\infty}^{+\infty}x[n]\e^{-\ju\omega n}
        \tag{The Analysis Equation}
    \end{align*}
    where the function $X(\e^{\ju\omega})$ is referred to as
    the \emph{discrete-time Fourier transform} or the \emph{spectrum} of $x[n]$.
\end{definition}

Since $\e^{\ju\omega}$ has a period of $2\pi$,
$X(\e^{\ju\omega})$ is always periodic with a period $2\pi$.
$\omega=0$ and $\omega=2\pi$ yield the same signal.
Therefore, usually we only consider the spectrum from $\omega=-\pi$ to $\omega=\pi$,
and signals at frequencies near $\omega=0$ or $\pm 2\pi$ can be considered as
\textbf{low-frequency}, while signals at frequencies near $\omega=\pm\pi$
can be considered as \textbf{high-frequency}.

\section{Fourier Transform for Periodic Sequences}

\begin{theorem}
    For a periodic sequence $x[n]$ whose FS coefficients are $\{a_k\}$,
    the Fourier transform $X(\e^{\ju\omega})$ is given by
    \begin{equation*}
        X(\e^{\ju\omega}) = \sum_{k=-\infty}^{\infty} 
        2\pi a_k \delta(\omega-k\omega_0)
    \end{equation*}
\end{theorem}
\begin{proof}
\begin{align*}
    X(\e^{\ju\omega})
    \quad\xrightarrow{\mathcal{F}^{-1}}\quad
    \frac{1}{2\pi}\int_{2\pi}
    X(\e^{\ju\omega})\e^{\ju\omega n}\dd{\omega}
    &=  \frac{1}{2\pi}\int_{W}^{W+2\pi}
    \sum_{k=-\infty}^{\infty}
    2\pi a_k \delta(\omega-k\omega_0)
    \e^{\ju\omega n}\dd{\omega} \\
    &= \sum_{k=-\infty}^{\infty}
    a_k \int_{W}^{W+2\pi}
    \delta(\omega-k\frac{2\pi}{N})
    \e^{\ju k\omega_0 n}\dd{\omega} \\
    &= \sum_{k=-\infty}^{\infty}
    a_k \e^{\ju k\omega_0 n}
    \left(\mycases{
        1 \qfor W \leqslant k \dfrac{2\pi}{N} < W+2\pi \\
        0 \qotherwise
    }\right) \\
    &= \sum_{k=\langle N\rangle}
    a_k\e^{\ju k\omega_0 n} \\
    &= x[n]
\end{align*}
\end{proof}



\section{Properties of CT Fourier Transform}

\subsection{}
Suppose we have:
\begin{gather*}
    \mycases[0.7]{
        x[n] \xrightleftharpoons[\mathcal{F}^{-1}]{\mathcal{F}}
        X(\e^{\ju\omega}) \\
        y[n] \xrightleftharpoons[\mathcal{F}^{-1}]{\mathcal{F}}
        Y(\e^{\ju\omega})
    }
\end{gather*}
Then we have the following results:
\begin{align*}
\shortintertext{Linearity}
    a\cdot x[n] + b\cdot y[n]
    \quad&\xrightleftharpoons[\mathcal{F}^{-1}]{\mathcal{F}}\quad
    a\cdot X(\e^{\ju\omega}) + b\cdot Y(\e^{\ju\omega})
\shortintertext{Time Shifting}
    x[n-n_0]
    \quad&\xrightleftharpoons[\mathcal{F}^{-1}]{\mathcal{F}}\quad
    \e^{-\ju\omega n_0} X(\e^{\ju\omega})
\shortintertext{Frequency Shifting}
    \e^{\ju\omega_0 n} x[n]
    \quad&\xrightleftharpoons[\mathcal{F}^{-1}]{\mathcal{F}}\quad
    X(\e^{\ju(\omega-\omega_0)})
\shortintertext{Time Reversal}
    x[-n]
    \quad&\xrightleftharpoons[\mathcal{F}^{-1}]{\mathcal{F}}\quad
    X(\e^{-\ju\omega})
\shortintertext{Scaling}
    x_{(k)}[n] =
    \mycases{
        x[\frac{n}{k}] & \text{for } \frac{n}{k}\in\mathbb{Z} \\
        0 & \text{otherwise}
    }
    \quad&\xrightleftharpoons[\mathcal{F}^{-1}]{\mathcal{F}}\quad
    X(\e^{\ju k\omega})
\shortintertext{Conjugation}
    x^\ast[n]
    \quad&\xrightleftharpoons[\mathcal{F}^{-1}]{\mathcal{F}}\quad
    X^\ast(\e^{-\ju\omega})
\shortintertext{Convolution}
    x[n]\ast y[n]
    \quad&\xrightleftharpoons[\mathcal{F}^{-1}]{\mathcal{F}}\quad
    X(\e^{\ju\omega}) Y(\e^{\ju\omega})
\shortintertext{Multiplication}
    x[n] y[n]
    \quad&\xrightleftharpoons[\mathcal{F}^{-1}]{\mathcal{F}}\quad
    \frac{1}{2\pi}\int_{2\pi}
    X(\e^{\ju\theta})Y(\e^{\ju\omega-\ju\theta})\dd{\theta}
\shortintertext{Differencing in Time}
    x[n] - x[n-1]
    \quad&\xrightleftharpoons[\mathcal{F}^{-1}]{\mathcal{F}}\quad
    (1-\e^{-\ju\omega}) X(\e^{\ju\omega})
\shortintertext{Accumulation in Time}
    \sum_{k=-\infty}^n x[k] = x[n] \ast u[n]
    \quad&\xrightleftharpoons[\mathcal{F}^{-1}]{\mathcal{F}}\quad
    \frac{1}{1-\e^{-\ju\omega}} X(\e^{\ju\omega})
    + \pi X(\e^{\ju 0}) \sum_{k=-\infty}^{\infty}\delta[\omega-2\pi k]
\shortintertext{Differentiation in Frequency}
    n\cdot x[n]
    \quad&\xrightleftharpoons[\mathcal{F}^{-1}]{\mathcal{F}}\quad
    \ju\dv{\omega} X(\e^{\ju\omega})
\shortintertext{Even-Odd Decomposition for Real Signals}
    x[n]\qq*{ real:} \Ev{x[n]}
    \quad&\xrightleftharpoons[\mathcal{F}^{-1}]{\mathcal{F}}\quad
    \Re{X(\e^{\ju\omega})} \\
    x[n]\qq*{ real:} \Od{x[n]}
    \quad&\xrightleftharpoons[\mathcal{F}^{-1}]{\mathcal{F}}\quad
    \ju\Im{X(\e^{\ju\omega})}
\shortintertext{Parseval's Relation for Aperiodic Signals}
    \int_{-\infty}^{\infty} \abs{x[n]}^2 \dd{n}
    \quad&=\quad
    \frac{1}{2\pi}
    \int_{2\pi} \abs{X(\e^{\ju\omega})}^2 \dd{\omega}
\end{align*}

\subsection{Symmetry}

\begin{itemize}
    \item $x[n]$ real $\quad\iff\quad$ $X(\e^{\ju\omega})=X^\ast(\e^{-\ju\omega})$
    \item $x[n]$ purely imaginary $\quad\iff\quad$ $X(\e^{\ju\omega})=-X^\ast(\e^{-\ju\omega})$
\end{itemize}

\paragraph{Conjugate Symmetry}
For real $x[n]$, its FT $X(\e^{\ju\omega})$ satisfies
\begin{itemize}
    \item $X(\e^{\ju\omega}) = X^\ast(\e^{-\ju\omega})$
    \item $\Re{X(\e^{\ju\omega})} = \Re{X(\e^{-\ju\omega})}$
    \item $\Im{X(\e^{\ju\omega})} = -\Im{X(\e^{-\ju\omega})}$
    \item $\abs{X(\e^{\ju\omega})} = \abs{X(\e^{-\ju\omega})}$
    \item $\arg X(\e^{\ju\omega}) = -\arg X(\e^{-\ju\omega})$
\end{itemize}

\paragraph{Parity}
\begin{itemize}
    \item $x[n]$ even $\quad\iff\quad$ $X(\e^{\ju\omega})$ even
    \item $x[n]$ odd $\quad\iff\quad$ $X(\e^{\ju\omega})$ odd
    \item $x[n]$ real \& even $\quad\implies\quad$
        $X(\e^{\ju\omega})$ real \& even
    \item $x[n]$ real \& odd $\quad\implies\quad$
        $X(\e^{\ju\omega})$ purely imaginary \& odd
\end{itemize}



\section{Basic Fourier Transform Pairs}

\begin{equation*}
    \sinc(\theta) := \frac{\sin(\pi\theta)}{\pi\theta}
\end{equation*}

\subsection{Table}

\begin{center}
\begin{tabular}
    {c c c}
    Signal & Fourier Transform & Fourier Series Coefficients \\
    \hline $\displaystyle
    \sum_{k=\langle N\rangle}
    a_k\e^{\ju k\omega_0 n}
    $ & $\displaystyle
    2\pi \sum_{k=-\infty}^{\infty} 
    a_k \delta(\omega-k\omega_0)
    $ & $
    a_k
    $ \\\hline $
    \e^{\ju\omega_0 n}
    $ & $\displaystyle
    2\pi \sum_{l=-\infty}^{\infty} \delta(\omega-\omega_0-2\pi l)
    $ & $
    \begin{matrix*}[l]
        \qif* \frac{\omega_0}{2\pi} = \frac{m}{N},
        & a_k = \mycases{
            1 & k={\scriptstyle m, m\pm N, m\pm 2N, \cdots} \\
            0 &\qotherwise*
        } \\
        \qif* \frac{\omega_0}{2\pi} \notin \mathbb{Q},
        & a_k \qq{is aperiodic}
    \end{matrix*}
    $ \\\hline $
    \cos(\omega_0 n)
    $ & $\displaystyle \!\!\!\!\!\!\!\!\!\!\!\!\!\!\!\!\!\!\!\!\!\!\!\!\!\!\!\!
    \pi \sum_{l=-\infty}^{\infty}
    \delta(\omega-\omega_0-2\pi l) + \delta(\omega+\omega_0-2\pi l)
    $ & $
    % \mycases{
    %     a_1 = a_{-1} = \frac{1}{2}\\ a_k = 0 \qotherwise
    % }
    $ \\\hline $
    \sin(\omega_0 n)
    $ & $\displaystyle \!\!\!\!\!\!\!\!\!\!\!\!\!\!\!\!\!\!\!\!\!\!\!\!\!\!\!\!
    \frac{\pi}{\ju} \sum_{l=-\infty}^{\infty}
    \delta(\omega-\omega_0-2\pi l) - \delta(\omega+\omega_0-2\pi l)
    $ & $
    % \mycases{
    %     a_1 = a_{-1} = \frac{1}{2\ju}\\ a_k = 0 \qotherwise
    % }
    $ \\\hline $
    \begin{matrix}
        \text{Periodic Square Wave} \\
        \mycases{
            1 & \abs{n} \leqslant N_1 \\
            0 & N_1 < \abs{n} \leqslant N/2 \\
        } \\
        \text{ with period of } N
    \end{matrix}
    $ & $\displaystyle
    2\pi\sum_{k=-\infty}^{\infty} a_k\delta(\omega-k\omega_0)
    $ & $\displaystyle
    a_k = \mycases{
        \dfrac{(2 N_1+1) \sin(k\omega_0)}{2 N\sin(k\omega_0/2)}
        & \qotherwise* \vspace{5pt}\\
        \dfrac{2 N_1 + 1}{N}
        & k = {\scriptstyle 0, \pm N, \pm 2  N, \cdots}
    }
    $ \\\hline $\displaystyle
    \sum_{k=-\infty}^{+\infty}\delta[n-k N]
    $ & $\displaystyle
    \frac{2\pi}{N}\sum_{k=-\infty}^{+\infty}\delta\qty(\omega-\frac{2\pi k}{N})
    $ & $\displaystyle
    a_k = \frac{1}{N}
    $ \\\hline $\displaystyle
    \mycases{
        1 & \abs{n} \leqslant N_1 \\
        0 & N_1 < \abs{n}
    }
    $ & $\displaystyle
    \frac{\sin[\omega(N_1+1/2)]}{\sin(\omega/2)}
    $ & \\\hline $\displaystyle
    \begin{matrix}
        \dfrac{\sin W n}{\pi n} \\
        0 < W < \pi
    \end{matrix}
    $ & $
    \begin{matrix}
        \mycases{
            1 & \abs{\omega} \leqslant W \\
            0 & W < \abs{\omega} \leqslant \pi
        } \\
        \text{with period of } 2\pi
    \end{matrix}
    $ & \\\hline $
    1
    $ & $\displaystyle
    2\pi\sum_{l=-\infty}^{\infty}\delta(\omega-2\pi l)
    $ & $
    a_k = \mycases{
        1 & k = 0, \pm N, \pm 2 N, \cdots \\
        0 &\qotherwise*
    }
    $ \\\hline $
    \delta[n]
    $ & $
    1
    $ & \phantom{$\mycases{1\\1}$} \\\hline $
    u[n]
    $ & $\displaystyle
    \frac{1}{1-\e^{-\ju\omega}} + \pi\sum_{l=-\infty}^{\infty}\delta(\omega-2\pi l)
    $ & \phantom{$\mycases{1\\1}$} \\\hline $
    \delta[n-n_0]
    $ & $
    \e^{-\ju\omega n_0}
    $ & \phantom{$\mycases{1\\1}$} \\\hline $
    a^{\abs{n} }\qq{with} \abs{a}<1
    $ & $\displaystyle
    \frac{1}{1-a\e^{-\ju\omega}} + \frac{1}{1-a\e^{\ju\omega}} - 1
    $ & \phantom{$\mycases{1\\1}$} \\\hline $
    a^n u[n] \qq{with} \abs{a}<1
    $ & $\displaystyle
    \frac{1}{1-a\e^{-\ju\omega}}
    $ & \phantom{$\mycases{1\\1}$} \\\hline $
    (n+1)a^n u[n] \qq{with} \abs{a}<1
    $ & $\displaystyle
    \frac{1}{(1-a\e^{-\ju\omega})^2}
    $ & \phantom{$\mycases{1\\1}$} \\\hline $\displaystyle
    \binom{n+r-1}{n}a^n u[n] \qq{with} \abs{a}<1
    $ & $\displaystyle
    \frac{1}{(1-a\e^{-\ju\omega})^r}
    $ & \phantom{$\mycases{1\\1}$} \\\hline
\end{tabular}
\end{center}

\subsection{Some Useful Techniques}

\begin{gather*}
    \delta[n] = \frac{\sin(\pi n)}{\pi n} = \sinc(n)
\end{gather*}

Expand the delta function to complex domain:
\begin{equation*}
    \delta(z) := \delta\big(\Re{z}\big)\delta\big(\Im{z}\big)
    \qfor z \in \mathbb{C}
\end{equation*}
Then we have:
\begin{equation*}
    \sum_{l=-\infty}^{\infty}\delta(\omega-2\pi l) = \delta(\e^{\ju\omega}-1)
    \qquad\qquad
    \sum_{l=-\infty}^{\infty}\delta(\omega+\omega_0-2\pi l)
    = \delta(\e^{\ju(\omega+\omega_0)}-1)
    \qquad\qquad
    \sum_{k=-\infty}^{+\infty}\delta\qty(\omega-\frac{2\pi k}{N})
    = \delta(\e^{\ju N\omega}-1)
\end{equation*}
\begin{align*}
    \e^{\ju\omega_0 n}
    \quad&\xrightleftharpoons[\mathcal{F}^{-1}]{\mathcal{F}}\quad
    2\pi \delta(\e^{\ju(\omega-\omega_0)}-1) \\
    \cos(\omega_0 n)
    \quad&\xrightleftharpoons[\mathcal{F}^{-1}]{\mathcal{F}}\quad
    \pi \left[
        \delta(\e^{\ju(\omega-\omega_0)}-1) +
        \delta(\e^{\ju(\omega+\omega_0)}-1)
    \right] \\
    \sin(\omega_0 n)
    \quad&\xrightleftharpoons[\mathcal{F}^{-1}]{\mathcal{F}}\quad
    \frac{\pi}{\ju} \left[
        \delta(\e^{\ju(\omega-\omega_0)}-1) -
        \delta(\e^{\ju(\omega+\omega_0)}-1)
    \right] \\
    \sum_{k=-\infty}^{+\infty}\delta[n-k N]
    \quad&\xrightleftharpoons[\mathcal{F}^{-1}]{\mathcal{F}}\quad
    \frac{2\pi}{N}\delta(\e^{\ju N\omega}-1) \\
    1
    \quad&\xrightleftharpoons[\mathcal{F}^{-1}]{\mathcal{F}}\quad
    2\pi \delta(\e^{\ju\omega}-1) \\
    u[n]
    \quad&\xrightleftharpoons[\mathcal{F}^{-1}]{\mathcal{F}}\quad
    \frac{1}{1-\e^{-\ju\omega}} + \pi\delta(\e^{\ju\omega}-1) \\
\end{align*}
\begin{equation*}
    a_k = \frac{X(\e^{\ju k\omega_0})}{2\pi\delta(\e^{\ju 2\pi k}-1)}
    \qquad\xrightleftharpoons[\mathcal{FS}]{\mathcal{FS}^{-1}}\qquad
    x[n]
    \qquad\xrightleftharpoons[\mathcal{F}^{-1}]{\mathcal{F}}\qquad
    X(\e^{\ju\omega}) = 2\pi a_{\omega/\omega_0}\delta(\e^{\ju N\omega}-1)
\end{equation*}

\section{Systems Characterized By LCCDEs}

An LTI system with initial rest characterized by the LCCDE
\begin{equation*}
    \sum_{k=0}^{N} a_k y[n-k] =
    \sum_{k=0}^{M} b_k x[n-k]
\end{equation*}
can be represented in the frequency domain as
\begin{equation*}
    H(\e^{\ju\omega}) =
    \frac{Y(\e^{\ju\omega})}{X(\e^{\ju\omega})} =
    \frac{\displaystyle\sum_{k=0}^{M} b_k \e^{-\ju\omega k}}
    {\displaystyle\sum_{k=0}^{N} a_k \e^{-\ju\omega k}}
\end{equation*}


\end{document}